\documentclass{article}
\usepackage[utf8]{inputenc}
\usepackage{geometry}
\geometry{a4paper, margin=1in}

\title{A Society of Learning Networks}
\author{Albert Jan van Hoek}
\date{\today}

\begin{document}

\maketitle

I have a vision of a society that finally understands what it is.

A society that sees humans not merely as competitors or consumers, but as nodes in a vast learning network stretching across generations.

\section*{The Nature of Mistakes and Feedback}

In that society, mistakes are not sources of shame but signals for improvement.

Feedback is not feared but welcomed.

Collaboration becomes the natural way forward, because learning flows best when minds connect.

\section*{The Flow of Knowledge}

We inherit the knowledge of those before us, we add what we discover, and we pass the baton onward.

When we recognize this, something simple becomes clear:

\textbf{Long-term collaboration is not an idealistic dream.}

It is simply what learning systems do when they understand themselves.

\section*{The Real Question}

The question is not whether we can build such systems.

The question is whether we finally recognize that we already are one.

\end{document}
